\chapter[Architecture: Habitable Sculpture?]{Is Architecture Just Sculpture You Can Live In?}
\section{An Uncooperative Door Handle}
Start with something small: the door handle you touch dozens of times every day.

Not the whole building, just the handle. Recall whether this happened last week: you approached a glass door and pulled. Nothing. You pulled harder. Still nothing. Then you noticed the sign: ``Push.'' Slightly embarrassed, you glanced around to see whether anyone noticed, pushed, and the door opened.

That small embarrassment contains a large question. American design researcher Donald Norman gave such doors the half-joking name ``Norman doors'': doors whose push-or-pull instructions you can never work out. His observation was that the mistake belongs to the door, not you. A flat metal plate says through its shape, ``Push me.'' A vertical handle says, ``Pull me.'' Put a pulling handle on a push-only door and design lies through its shape; you merely believed it.

Notice what happened. Metal and glass, without mouths or words---the push sign is the designer's later patch---direct your body. They tell you where to put your hand, how much force to use, and which way to move. Expand those instructions to the building: stair width sets how many can walk abreast, corridor length how many seconds from classroom to toilet, door height whether you duck, and a square's orientation where sunlight falls on you.

Architecture never merely stands there. It is instructions written for your body. Who writes them, and for whom? Those are this chapter's questions.

\section{Entering Hagia Sophia}
Come with me into a scene.

Time: an early winter morning in 562 CE. The building was completed in 537 and its earthquake-damaged dome has just been rebuilt. Place: Constantinople, today's Istanbul in Turkey. You follow a crowd toward the immense cathedral of Hagia Sophia.

You tilt your head long before reaching it, not from piety but because your neck responds on its own. The dome stands over fifty meters high, more than a dozen stories; in the sixth century, the tallest thing you have ever seen is a three-story tower. The doorway is disproportionately tall. Crossing its threshold, deliberately set a little below the outside ground, your body understands before your mind: this place is outside ordinary life.

Inside there are no lamps, at least none familiar to you. Light pours through a dense ring of little windows at the dome's base, making it seem supported by light rather than columns, floating in midair. Frankincense scents the cool, echoing air. Someone coughs ahead and the sound rolls among stones for seconds. Marble covers the walls, deliberately arranged: masons split slabs and open them like books, matching their patterns into enormous flowers. You have heard these stones came from quarries across the empire, some shipped hundreds of kilometers by sea.

Behind you, a rural pilgrim wipes her eyes. Later she tells others that heaven truly seemed closer to earth that moment.

What is the building doing? Housing people? Hardly: thousands may stand inside, but nobody sleeps or eats there. Is it sculpture? It exceeds any sculpture in size, and you must enter to experience it; sculpture does not let you walk inside. It does a third thing: stone, light, height, and echo rearrange body and mind, making you look up, tread softly, feel small, and feel happily small.

Nine centuries later, in 1453, the Ottomans captured the city and converted Hagia Sophia into a mosque. Plaster covered Christian mosaics, slender minarets rose around it, and Arabic calligraphic roundels hung inside. In 1935, the new Turkish republic made it a museum and uncovered the mosaics. In 2020, it became a mosque again. The same stones and dome changed identity four times in fifteen centuries. The building scarcely changed, but each generation asked a different question: whose are you now?

\section{Sculpture or Instrument?}
Set out the disagreement before answering it.

``Is architecture just habitable sculpture?'' has two very different groups behind it.

One says architecture is certainly art, and a high art. A widespread European saying, often attributed to Goethe, calls it frozen music: rhythm, proportion, and emotion solidified in stone. Buildings, sculptures, and symphonies should therefore be judged with the same language: are they beautiful and moving? Hagia Sophia's dome, the Parthenon's colonnades, and framed garden views in Suzhou are first aesthetic objects.

The other says no. Sculpture need not shelter from rain, direct people to toilets, or avoid killing residents when it collapses; buildings must. Architecture is first an instrument, an enormous chair or clothing worn by society. A chair is judged first by comfort, not resemblance to a symphony. Only outsiders viewing photographs mistake a building for sculpture; its inhabitants do not.

Both sides have powerful evidence. Almost no civilization settles for mere adequacy: early huts had carved lintels, and Gothic spires rose beyond any practical value while increasing collapse risk. Pure utility cannot explain the extras. Yet even sacred buildings must drain rain, keep warm, and admit thousands. Most historical buildings---homes, warehouses, bridges, aqueducts---were never considered art, yet constitute ninety-nine percent of human construction.

Perhaps both questions are too narrow. Recall the handle and Hagia Sophia's threshold. Architecture's deepest action is neither looking good nor working well, but \textbf{arranging}: where bodies move, pause, meet, and remain unseen. Neither side has addressed this dimension. When it enters later, you will discover that disputes over architecture are never only architectural.

\section{Who Is Invited In, Who Is Kept Out?}
A building's most revealing feature is not its roof but its threshold: who may cross and who may not.

Consider ancient Egypt's Karnak temple at modern Luxor, among history's largest religious complexes. Its grandeur opened to very few. The design narrows and darkens inward: monumental pylons and outer courts allow festival crowds; a hall of hundreds of columns admits fewer; progressively smaller, lower rooms lead to the tiny dark inner sanctuary housing the deity, normally open only to pharaoh and the highest priests. Space is a hierarchy written in stone. Distance from the god depends not on your heart but your status. Each doorway reconfirms who you are.

Now a typical Beijing courtyard house, or siheyuan. The gate lies southeast, with a screen wall blocking outside sightlines. Servants and the master's reception rooms occupy the outer courtyard; women and children live farther inside. The north main building faces south and receives the best light, reserved for elders; younger generations occupy east and west wings. Bricks translate seniority and the separation of inner and outer into daily bodily experience: which door you leave, whether sunlight reaches you, whether you see the street. A child need hear no lecture; walking has already taught the household hierarchy.

Palaces scale this up to the state. Beijing's Forbidden City unfolds along a north--south axis with major halls on it. Nearer the center, more gates, emptier courtyards, and higher steps make a minister entering by a side gate cross successive spaces and climb platforms. He is not merely walking; space repeatedly tells him his distance from power. Louis XIV went further at Versailles, placing his bedroom centrally and making waking and retiring ceremonies. Nobles entered by rank, their rooms and positions dictated by the king. A palace's sharpest tactic is not exclusion but \textbf{admitting you, then making you feel arranged at every turn}.

Gardens have another intention. Suzhou's private gardens often belonged to retired or frustrated scholars. Small plots deliberately resist a single comprehensive view: winding corridors, framed borrowed scenery, and a pond suddenly beyond a white wall. Concealment and revelation enlarge space through movement, arranging rhythm and mood rather than hierarchy, a manageable little world for its owner. Japanese dry-landscape gardens simplify further: white sand and a few rocks without water, inviting prolonged contemplation from a veranda. Even wandering disappears; space arranges only sitting and looking. Arrangement can be a gift to oneself, not necessarily oppression.

A typical mosque offers another answer. Its broadest area is often a large courtyard with water, theoretically open to all believers. The prayer hall has no divine image, seating, or progressively raised altar, only a recessed mihrab marking Mecca's direction. Worshippers stand shoulder to shoulder in rows, rich and poor, old and young, facing together. The space declares equality before God, physically made through a level floor and close ranks. That is the tradition's claim, however. Women's prayer areas often occupy side rooms or upper floors, and differences of position and sightlines remain debated within communities. No tradition speaks with only one voice.

Gothic cathedrals differed. Medieval Europe's great churches were among their cities' most accessible large buildings: they could host markets, shelter refugees, and let poor visitors view nave stained glass without payment. Those pictures became known as the poor person's Bible, stories for people unable to read. Remember this corrective: grand religious buildings do not always exclude; they can be exceptionally inclusive public spaces. Assess thresholds building by building.

Homes are the planet's most numerous buildings and least discussed in art history. From northern Shaanxi cave dwellings and Fujian tulou to nomadic felt tents and Japanese machiya, they are rarely built for display, yet honestly record daily life. A tulou circles dozens of households into one defensive unit behind a closed gate, revealing former threats. A felt tent dismantled and loaded within an hour reveals a life following pasture.

Monuments expose most directly whom architecture represents. Many imagine Egyptian pyramids built by thousands of enslaved people under whips. Recent excavations of workers' settlements, bakeries, graves, and medical traces suggest to mainstream Egyptologists organized labor teams receiving rations, often serving in rotation, with some proud enough to be buried nearby. Grandeur does not automatically equal oppression. Yet not enslaved does not mean easy or voluntary labor: royal tombs still demanded enormous state-organized human effort. Stone names the commemorated pharaoh and his eternity; the unnamed workers' graves reveal those left silent.

\section[Five Questions for Reading Buildings]{Thinking Bridge: Five Questions for Reading a Building}
Can we read any building, from school to ancient temple, through fixed questions rather than inspiration? Yes. Architectural analysis can be unpacked into five questions that largely open a building to you.

\textbf{First: function---what must happen here?}

Not whether it is a school or temple, but which actions its occupants need: sitting, standing, lying, gathering, dispersing, seeing, being seen, hiding, waiting. A hospital waiting room must let anxious people wait without falling apart. Chair arrangements and queue displays reveal how well it succeeds.

\textbf{Second: structure---what keeps it standing?}

Load-bearing walls or columns? Beams, arches, or frames? Structure sets spatial freedom. Greek stone beams span little, requiring forests of columns. Roman concrete arches enlarged spans tenfold and opened huge interiors. Reinforced-concrete frames free walls from supporting loads, allowing all-glass walls. Structure is not merely engineering; it is the constitution governing what space can become.

\textbf{Third: materials---what, where, and whose hands?}

Materials are never purely technical. Wood, stone, brick, concrete, and glass have different qualities, costs, and transport ranges. A local rammed-earth home and a temple of imported marble speak different political languages: we live in this soil, versus our reach extends hundreds of kilometers.

\textbf{Fourth: space---how does the body move?}

The most important question. Do not stand outside looking at photographs; follow an imagined body. Which entrance, first view, necessary turn or pause, overlook, or exposed position? The walk reveals spatial grammar. Try your school: why classrooms on both corridor sides, teachers' offices at stairways, and a playground visible from every classroom window?

\textbf{Fifth: symbolism---what feelings are intended?}

Height makes you look up, size makes you feel small, symmetry suggests order, and dark winding routes slow your steps. This is not mysticism but stable responses between body and space that work worldwide.

The questions interlock: function sets needed space, structure possible space, materials its feel, and symbolism often grows from the others. Next time a building impresses you, move beyond admiration. Grandeur is an engineering achievement that can be analyzed.

\section{Two Plans from Two Thousand Years Ago}

Read some original material. People more than two millennia ago already described architecture's tasks, with surprisingly similar Eastern and Western answers.

China's Kaogong ji, or Record of Trades, in the Rites of Zhou records state-craft standards around the Warring States period. Its ideal capital has this plan:

\begin{quote}
The craftsman lays out the capital as a square nine li on each side, with three gates on each side. Within are nine north--south and nine east--west avenues; the north--south roads are nine carriage tracks wide.
\end{quote}
(Rites of Zhou, Record of Trades)

In brief: a square capital nine li per side, three gates per wall, nine avenues in each direction, and north--south roads wide enough for nine carriages. A few dozen characters arrange an entire city: square outline, strict grid, and road ranks measured by wheel tracks. This is not the voice of an engineering manual but an \textbf{ideal}, describing what a ritually fitting capital ought to be. Later Chang'an and Beijing echo this grid. Architecture here means order: heaven is not round, but earth is square; earth is not chaotic because the city has its coordinates.

In the West, the first-century-BCE Roman architect Vitruvius wrote Ten Books on Architecture, the only complete surviving ancient Western architectural treatise. Good buildings require firmitas, utilitas, and venustas: \textbf{strength, utility, and beauty}. They must stand, function, and please the eye; none is optional.

Compare them. Beauty has an explicit place in Vitruvius's triad. Romans answered the art-versus-instrument dispute two thousand years ago: both sides are right, neither complete. The Record of Trades concerns not one building's beauty or strength but a whole city's arrangement. Its crucial property is \textbf{ritual propriety}: roads, gates, and dimensions embody political rank.

Together they open a field far larger than habitable sculpture. Architecture answers to bodies through strength and utility, eyes through beauty, and order through propriety and rules. The latter deserves deeper scrutiny, because order always has authors and costs.

\section{Why Are Cathedrals So Tall?}
Practice comparing explanations. From the twelfth century, French, English, and German cities competed furiously to build ever-taller cathedrals. Beauvais reached a forty-eight-meter vault, collapsed during construction, was rebuilt, collapsed again, and remains unfinished---yet people would not lower it. Why?

Separate four things: taller buildings and collapses are events; why they happened concerns competing explanations; glory to God in documents is participants' self-understanding; your judgment comes last.

\textbf{Explanation one: theology---height approaches God.}

Participants said height, light, and rising lines directed believers toward heaven. Bundled columns, pointed arches, and ribs pull the eye upward; stained glass makes otherworldly red and blue light. This explanation has solid grounds: funders and builders said it, and design serves the experience. Its limit is treating their self-understanding as cause. Glorifying God mattered in the eleventh century too; why the twelfth-century leap? Desire alone cannot explain timing.

\textbf{Explanation two: technology---height became possible.}

Pointed arches, flying buttresses, and ribbed vaults formed a new structural toolkit, conducting roof loads outward along curves so walls could thin, gain windows, and rise. This explains timing: tools arrived with ambition. But technology explains can, not want. Inventing a shovel does not automatically inspire the world's deepest hole. Why Beauvais insisted on height despite collapse is beyond the shovel's answer.

\textbf{Explanation three: civic competition---height addresses neighbors.}

Historians note the fiercest competition in commercially advanced northern France. Semiautonomous urban communities of merchants and craft guilds had practical pride: a cathedral advertised the city, attracting pilgrims, traders, and markets, strengthening its position against neighbors, bishops, and kings. Dedicated to God, it also displayed achievement to rivals. Competition explains successive imitation, but reducing everything to advantage cannot explain lifelong savings donated after repeated collapse. Cold accounts should not contain such devotion.

The explanations grasp desire, ability, and interests, not mutually exclusive rivals but parts of one machine. Without desire, competition lacks meaning; without capacity, it cannot begin; without rivalry, neither is pushed toward rebuilding ever higher. Apply this elsewhere: skyscrapers, dams, theme parks. Find desire, ability, and competition, and apparent madness becomes intelligible.

\section[The Panopticon and Spatial Power]{Deeper Challenge: The Panopticon and Spatial Power}
Now introduce the chapter's weightiest theory, making architecture's arrangement of social relationships unsettling.

In 1791, Jeremy Bentham published the Panopticon prison design. Terrifyingly simple: a ring of cells around a central watchtower. Each cell receives outside light and opens toward the tower. A guard can turn to inspect anyone, while backlit prisoners cannot see inside the tower and \textbf{never know whether they are currently watched}. Few guards are needed because possible observation does the work. Bentham expected prisoners to regulate themselves: power's whip moves into their minds. His complete design was not built in his lifetime, but the plan became a major twentieth-century intellectual device.

Michel Foucault revived it in Discipline and Punish in 1975; the following paraphrases his argument. The design is not merely a prison but a diagram of modern society. Traditional power beats bodies through expensive, ugly torture and chains; modern power supervises through respectable, cheaper schedules, records, examinations, rankings, cubicles, and cameras. Architecture silently assists. Schools sort children by age into identical rooms, desks facing one direction, sometimes with inspection windows. Hospitals sort diseases across floors, carefully space beds, and place nurses' stations for visibility. Open-plan offices expose screens to passing eyes. Such spaces perform surveillance without a particular guard.

The theory is sharp, but weigh its limits. Read every classroom as prison and it starts devouring everything. The same arrangement has taught billions to read over five centuries; hospital organization aids surveillance and swift emergency care. Arranging bodies is itself neutral. Ask each time what purpose it serves, who benefits, and whether occupants can refuse. A theory supplies glasses, not a ready verdict. Each building still needs fresh examination.

Keep one unanswered question: Bentham's prisoners knew they might be watched. When you enter shops, stations, and classrooms today, how often do you truly know, rather than your body simply having become accustomed? Does the camera-filled city exceed even the Panopticon, or exceed what that two-century-old plan can explain?

\section{Monuments, Benches, and Ramps Today}
\textbf{First, monuments.} Washington's Vietnam Veterans Memorial opened in 1982. Its competition-winning designer, Maya Lin, was a twenty-one-year-old Chinese American student. Instead of obelisk or heroic soldier, two black granite walls cut a shallow V into the ground and descend gradually, bearing over fifty-eight thousand dead people's names in chronological order. Walk downward and polished stone reflects your face over the names.

The proposal caused uproar. Some veterans considered blackness, descent, and absent heroes shameful, a wound in the earth; they wanted an upright monument. Organizers compromised by later adding a conventional bronze soldier statue nearby. Yet the completed wall became among America's most visited memorials, where people lingered, found fathers, brothers, and comrades, took pencil rubbings, and left letters, photographs, and medals. Many fierce critics tearfully withdrew their words.

The case exposes who is remembered and silenced. Traditional monuments honor sacrifice's meaning---flag, hero, nation---with names in the background. Lin's wall honors the people sacrificed, fifty-eight thousand particular names, leaving meaning to visitors. Neither satisfies everyone: one submerges individuals beneath grandeur, the other allegedly evades whether the war was worthwhile. Monument disputes concern \textbf{authority over the narrative}, who may define the sacrifice. Arguments over removing old statues or inscribing new walls still put that question on trial worldwide.

\textbf{Next, benches.} Nearby parks probably contain benches divided by extra armrests or separate single seats. Officially, armrests help older people stand, but also prevent lying down. Hostile architecture includes narrow sloping bus-stop seats, concrete spikes under overpasses, and nighttime shopfront sprinklers, all intended to keep homeless people from staying.

Steelman the other side seriously. Supporters are not cartoon villains: shops owe customers and staff safety and cleanliness, residents worry about children's routes home, municipalities face complaints, and homelessness originates in housing, healthcare, and support systems. Expecting one convenience store's missing sprinkler to solve it shifts society's responsibility onto a distant endpoint. Each reason has weight. Yet throughout this account, people sleeping outside appear as a problem to manage, not street users. The deepest problem is \textbf{defining public space from only one kind of person's viewpoint}. Benches are for people, with people silently narrowed.

\textbf{Finally, ramps.} A curb is one step for a walker but a wall for a wheelchair user. In the 1970s, disability-rights groups won curb ramps in some American cities, despite objections to spending everyone's money for a small minority. Their users exceeded expectations: parents with strollers, travelers with suitcases, scooter-riding children, delivery workers with carts, and older people. The curb-cut effect names how designs for the most disadvantaged often help everyone. Accessibility is not charity for special others but a reconsideration of normal. A route a fifteen-year-old manages with strong knees should also serve an eighty-year-old, a mother pushing a stroller, and a wheelchair user. Broken tactile paths, ramps ending at barriers, and perpetually locked reservation-only subway lifts reveal who a city assumes its citizens are.

At city scale, consider Brasilia, built on a plateau in 1960. Strict functional zones and symmetrical axes make a birdlike aerial plan, modernist planning's dream. Residents discovered immense crossings, no corner shops or wandering lanes between districts, and bleak walking. Around then, Jane Jacobs defended the opposite street in The Death and Life of Great American Cities; this paraphrases her argument. Good neighborhoods mix homes, shops, workshops, old and new buildings, and constantly occupied sidewalks. Neighbors' eyes provide security: eyes on the street. Order drawn from the sky versus order grown from sidewalks: the Record of Trades and daily necessities still struggle within one city two millennia later.

\section{Who Is Your City's Default Citizen?}
Return to the door handle.

Metal says push for a designer; corridor width says one at a time for a school; temple thresholds say no farther for pharaoh; a sunken wall says a nation owes these names an answer; an obstructed tactile path says a city forgot you. Architecture never lies; it speaks for its author. Who is that author, and whom is it addressing?

Do one thing, then answer one question.

Use the five questions on your school. Why do classrooms face one direction? How are sightlines between building and playground arranged? Which doors stay locked, and why? If your classmate came in a wheelchair, how many barriers lie between gate and desk? Write it down. You may see your familiar building for the first time.

The question: \textbf{who should be the default user when public space is designed?}

Choose the strongest, most standard user, young, able-bodied, and hurried, and construction is cheapest and space most efficient; that person does form the majority. Choose the least conveniently situated, using wheelchair, stroller, cane, or no sight, and everyone's passage costs slightly more, but nobody is barred; the curb-cut effect suggests eventual benefit to all. Both accounts work, but imply different meanings of public: belonging to the majority, or deserving the name only when minorities too are included?

Give your answer and reasons. Before deciding, recall Hagia Sophia's four identities in fifteen centuries, each assigned by those empowered to answer whom it served. Your city, school, and neighborhood bench answer the same question now. This time, you can hear what they are saying.
